% ---- T-SQL listing style (safe to keep if already defined in main.tex) ----
\lstdefinestyle{tsql}{
  language=SQL,
  basicstyle=\ttfamily\footnotesize,
  keywordstyle=\color{blue!70!black}\bfseries,
  commentstyle=\color{green!40!black}\itshape,
  stringstyle=\color{red!60!black},
  numbers=left, numberstyle=\tiny\color{gray},
  frame=single, breaklines=true, columns=fullflexible,
  showstringspaces=false, tabsize=2
}
\lstset{style=tsql}

\chapter{Query Optimization and Query Processing Analysis}
\label{chap:query-processing}

\section{Introduction}

This chapter examines how SQL Server processes the project's representative queries, with emphasis on query optimization, cardinality estimation, statistics, join algorithms, scalar calculations, and parallel execution. The analysis focuses on the baseline database configuration before applying the additional indexes.

The investigation uses actual execution plans, operator properties, SQL Server statistics, and runtime measurements. The execution plans provide evidence of the access paths and processing operators selected by the optimizer, while the statistics and cardinality estimates help explain its decisions. Runtime and I/O observations are used to examine the resulting execution behaviour.

% =============================================
% Q1: Patient and MRI Scan Retrieval
% =============================================

\section{Q1: Patient and MRI Scan Retrieval}

\subsection{Query Characteristics}

Q1 joins \texttt{Patients} and \texttt{MRI\_Scans} using \texttt{PatientID} and filters for patient \texttt{'0001'}. The query returns the patient's information together with the matching MRI scan records.

\subsection{Execution Plan and Join Selection}

Figure~\ref{fig:q1-plan} presents the baseline actual execution plan for Q1. The plan uses a Clustered Index Seek on \texttt{Patients}, followed by a Clustered Index Scan on \texttt{MRI\_Scans}. The two inputs are combined using a Nested Loops (Inner Join) operator.

\begin{figure}[htbp]
\centering
\planfig{image/role_3/q1_execution_plan.png}{0.5\textwidth}{Q1 baseline actual execution plan}
\caption{Baseline actual execution plan for Q1}
\label{fig:q1-plan}
\end{figure}

The seek on \texttt{Patients} returns one row, which forms a small outer input for the Nested Loops operator. The inner input scans \texttt{MRI\_Scans} because the baseline database does not have an index on its \texttt{PatientID} column. The scan reads all 48,345 rows and returns 65 matching rows.

The relevant properties of the Clustered Index Scan are summarized in Table~\ref{tab:q1-scan}.

\newpage

\begin{table}[htbp]
\centering
\caption{Q1 Clustered Index Scan properties for \texttt{MRI\_Scans}}
\label{tab:q1-scan}
\begin{tabular}{lr}
\toprule
Property & Value \\
\midrule
Actual rows read & 48,345 \\
Actual rows returned & 65 \\
Estimated rows returned & 79.875 \\
Estimated rows to be read & 48,345 \\
Actual executions & 1 \\
Estimated operator cost & 1.44387 (98\% of plan cost) \\
\bottomrule
\end{tabular}
\end{table}

Figure~\ref{fig:q1-operator} provides the detailed operator properties used to obtain these values.

\begin{figure}[htbp]
\centering
\planfig{image/role_3/q1_scan_properties.png}{0.5\textwidth}{Q1 MRI_Scans Clustered Index Scan properties}
\caption{Actual and estimated properties of the Q1 Clustered Index Scan}
\label{fig:q1-operator}
\end{figure}

The optimizer's choice of Nested Loops is consistent with the one-row outer input. However, the absence of an index on the inner input's join column means that the scan must examine the entire \texttt{MRI\_Scans} table. Therefore, the principal inefficiency in this baseline plan is the amount of data read, rather than a large cardinality-estimation error.

\subsection{Cardinality Estimation}

The Clustered Index Seek on \texttt{Patients} estimates one row and returns one row. For the scan on \texttt{MRI\_Scans}, SQL Server estimates 79.875 matching rows, while the actual number is 65.

The Q-error is calculated as:

\[
Q\text{-error}=
\max\left(
\frac{79.875}{65},
\frac{65}{79.875}
\right)
\approx 1.23.
\]

The estimate is therefore approximately 23\% higher than the actual result. This is a relatively small estimation error for this query.

\subsection{Statistics for \texttt{MRI\_Scans.PatientID}}

The following query retrieves the statistics information for \texttt{MRI\_Scans.PatientID}, including its header, density vector, and histogram.

\begin{lstlisting}[caption={Retrieving Q1 statistics for MRI\_Scans.PatientID},
                   label={lst:q1-statistics}]
DBCC SHOW_STATISTICS (
    'dbo.MRI_Scans',
    '_WA_Sys_00000002_4BAC3F29'
);
GO
\end{lstlisting}

The statistics metadata and density information, together with histogram entries, are summarized in Tables~\ref{tab:q1-statistics} and \ref{tab:q1_patientid_histogram}, respectively.

\begin{table}[htbp]
\centering
\caption{Statistics Metadata and Density Information for \texttt{MRI\_Scans.PatientID}}
\label{tab:q1-statistics}
\begin{tabular}{lr}
\toprule
Property & Value \\
\midrule
Rows & 48,345 \\
Rows sampled & 32,599 \\
Sampling percentage & 67.43\% \\
Steps & 200 \\
Density & 0.02050357 \\
Average key length & 4 \\
String index & YES \\
Unfiltered rows & 48,345 \\
Persisted sample percent & 0 \\
\bottomrule
\end{tabular}
\end{table}

\begin{table}[htbp]
\centering
\small
\caption{Histogram Entries for \texttt{MRI\_Scans.PatientID}}
\label{tab:q1_patientid_histogram}
\begin{tabular}{|r|r|r|r|r|}
\hline
RANGE\_HI\_KEY & RANGE\_ROWS & EQ\_ROWS &
DISTINCT\_RANGE\_ROWS & AVG\_RANGE\_ROWS \\
\hline
0001 & 0 & 79.87501 & 0 & 1 \\
\hline
0003 & 52.05838 & 105.0208 & 1 & 51.85342 \\
\hline
0007 & 182.948 & 144.9583 & 3 & 60.8224 \\
\hline
0009 & 74.36911 & 94.66667 & 1 & 74.07632 \\
\hline
0011 & 69.90696 & 103.5417 & 1 & 69.63174 \\
\hline
\multicolumn{5}{|c|}{$\vdots$} \\
\hline
0565 & 188.8975 & 35.5 & 2 & 94.16961 \\
\hline
0568 & 80.31863 & 109.4583 & 1 & 80.00242 \\
\hline
0571 & 239.4685 & 127.2083 & 2 & 119.3804 \\
\hline
0573 & 65.44482 & 221.875 & 1 & 65.18716 \\
\hline
0575 & 0.1880675 & 79.87501 & 0 & 95.16214 \\
\hline
\end{tabular}
\end{table}

The histogram contains an exact entry for \texttt{PatientID = '0001'}, with \texttt{EQ\_ROWS = 79.87501}. This corresponds to the optimizer's estimated cardinality of approximately 79.875 rows in the execution plan. Compared with the 65 actual matching rows, this represents an overestimation of approximately 23\%.

This indicates that the histogram provides a reasonably close estimate for the equality predicate. However, the statistics estimate does not change the access path: without an index on \texttt{MRI\_Scans.PatientID}, SQL Server still performs a Clustered Index Scan.

\subsection{Runtime and I/O Observations}

The captured execution plan reports 1,880 logical reads, 3 physical reads, and 1,874 read-ahead reads for the \texttt{MRI\_Scans} scan. The query consumed approximately 369 ms of CPU time, while its elapsed time was approximately 76.3 seconds.

The substantial difference between CPU and elapsed time indicates that the query spent considerably more time waiting than actively processing data. The most prominent reported wait was \texttt{SLEEP\_BPOOL\_STEAL}, accounting for approximately 75.9 seconds. However, the available execution-plan data does not establish the underlying cause of this wait, so the elapsed time should not be treated as a representative measure of the scan's processing cost without further investigation.

The plan also reports that it was retrieved from cache. This means the captured execution reused a cached execution plan rather than requiring a fresh compilation. The runtime and I/O measurements are therefore recorded as observations from this particular execution, rather than as definitive performance characteristics of Q1.

\subsection{Nested Loops versus Hash Match for Q1}

To test the optimizer's join choice, Q1 was also executed with \texttt{OPTION (HASH JOIN)}. Table~\ref{tab:q1-join-comparison} compares the two plans.

\begin{table}[htbp]
\centering
\caption{Q1: Nested Loops (chosen) versus Hash Match (forced)}
\label{tab:q1-join-comparison}
\begin{tabular}{lrr}
\toprule
Metric & Nested Loops & Hash Match \\
\midrule
Estimated subtree cost & 1.47069 & 1.4915 \\
Join operator cost & 0.0235 (2\%) & 0.0443 (3\%) \\
Memory grant (KB) & 0 & 1,072 \\
Degree of parallelism & 1 & 1 \\
Estimated / actual join rows & 79.875 / 65 & 79.875 / 65 \\
Logical reads, \texttt{MRI\_Scans} & 1,880 & 1,880 \\
\bottomrule
\end{tabular}
\end{table}

Both plans read the same 1,880 pages of \texttt{MRI\_Scans}, because the inner table must be scanned in full without an index on \texttt{PatientID}. The forced Hash Match raises the join operator cost from 0.0235 to 0.0443 and requires a 1,072~KB memory grant, with no spill to tempdb (zero Workfile and Worktable reads). The total estimated cost is therefore about 1.4\% higher than the optimizer's chosen plan, which supports its selection of Nested Loops.

However, the scan accounts for most of the plan cost in both cases, so the join algorithm has only a minor influence on Q1. The estimated cardinality of 79.875 rows is unchanged, so the forced join does not alter the Q-error of 1.23.

% =============================================
% Q2: Filtered Scan Retrieval by Pulse Sequence
% =============================================

\section{Q2: Filtered Scan Retrieval by Pulse Sequence}

\subsection{Query Characteristics}

Q2 retrieves the T2-weighted sagittal series of patient \texttt{'0001'} by joining \texttt{Patients} and \texttt{MRI\_Scans} on \texttt{PatientID} and applying the predicate \texttt{SeriesDescription LIKE '\%T2\%SAG\%'}. Clinically, radiologists inspect T2-weighted sagittal series to assess disc herniation and cord compression, so this filtered retrieval is a representative workload for the reporting console.

\subsection{Execution Plan and Join Selection}

Figure~\ref{fig:q2-plan} presents the baseline actual execution plan for Q2. As with Q1, the plan combines a Clustered Index Seek on \texttt{Patients} with a full Clustered Index Scan of \texttt{MRI\_Scans} through a Nested Loops (Inner Join) operator. The scan contributes 96\% of the total plan cost (0.85990 out of 0.89564).

\begin{figure}[htbp]
\centering
\planfig{image/role_3/q2_execution_plan.png}{0.5\textwidth}{Q2 baseline actual execution plan}
\caption{Baseline actual execution plan for Q2}
\label{fig:q2-plan}
\end{figure}

The leading wildcard in the \texttt{LIKE} predicate makes it unsargable: no index can serve a filter of the form \texttt{LIKE '\%T2\%SAG\%'}, so the predicate is evaluated as a residual filter on every row of the scan. The scan reads all 48,345 rows and returns the 15 rows of patient \texttt{'0001'} whose series description matches the pattern.

The relevant properties of the Clustered Index Scan are summarized in Table~\ref{tab:q2-scan}.

\begin{table}[htbp]
\centering
\caption{Q2 Clustered Index Scan properties for \texttt{MRI\_Scans}}
\label{tab:q2-scan}
\begin{tabular}{lr}
\toprule
Property & Value \\
\midrule
Actual rows read & 48,345 \\
Actual rows returned & 15 \\
Estimated rows returned & 13.8719 \\
Estimated rows to be read & 48,345 \\
Actual executions & 1 \\
Estimated operator cost & 0.85990 (96\% of plan cost) \\
\bottomrule
\end{tabular}
\end{table}

Figure~\ref{fig:q2-operator} provides the detailed operator properties used to obtain these values.

\begin{figure}[htbp]
\centering
\planfig{image/role_3/q2_scan_properties.png}{0.5\textwidth}{Q2 MRI_Scans Clustered Index Scan properties}
\caption{Actual and estimated properties of the Q2 Clustered Index Scan}
\label{fig:q2-operator}
\end{figure}

The predicate list shown in the scan properties is \texttt{(SeriesDescription LIKE '\%T2\%SAG\%') AND (@1 = PatientID)}. Both the \texttt{LIKE} expression and the \texttt{PatientID} equality appear as residual predicates rather than seek arguments, a direct consequence of the scan being the chosen access method in the absence of any index on the join column. The plan also reports a missing-index hint suggesting a nonclustered index on \texttt{MRI\_Scans.PatientID} with \texttt{SeriesDescription} and \texttt{ImagePath} as included columns, with an impact of 99.0316\% --- the largest improvement opportunity in the plan. The optimizer therefore identifies the missing join-column index, not a \texttt{SeriesDescription} index, as the dominant optimization target.

\subsection{Cardinality Estimation}

The Clustered Index Seek on \texttt{Patients} estimates one row and returns one row. For the scan on \texttt{MRI\_Scans}, SQL Server estimates 13.8719 matching rows, while the actual number is 15.

The Q-error is calculated as:

\[
Q\text{-error}=
\max\left(
\frac{13.8719}{15},
\frac{15}{13.8719}
\right)
\approx 1.08.
\]

The estimate is therefore approximately 8\% lower than the actual result, a small error that one would not expect for an unsargable \texttt{LIKE} predicate. The reason is that the optimizer does not fall back to a flat default selectivity: dividing the estimated rows by the \texttt{PatientID} estimate (13.8719 / 79.875) gives an implied \texttt{LIKE} selectivity of approximately 0.174, which is close to the global proportion of \texttt{MRI\_Scans} rows whose series description matches the pattern (9,809 out of 48,345, or 20.3\%). For patient \texttt{'0001'}, all 15 matching slices belong to the single most frequent series in the database, \texttt{T2\_TSE\_SAG\_384\_0002}, so the step-level histogram estimate coincides well with the true value. For patients whose matching slices lie in the long tail of the distribution, the same estimation method would be considerably less reliable.

\subsection{Statistics for \texttt{MRI\_Scans.SeriesDescription}}

The following query retrieves the statistics information for \texttt{MRI\_Scans.SeriesDescription}, including its header, density vector, and histogram. The statistics object \texttt{\_WA\_Sys\_00000004\_4CA06362} is auto-created on \texttt{SeriesDescription}; auto-created statistics names are generated per database instance, so the name should be verified on other deployments (see the lookup query in \texttt{scripts/sql/03\_role\_3\_support.sql}).

\begin{lstlisting}[caption={Retrieving Q2 statistics for MRI\_Scans.SeriesDescription},
                   label={lst:q2-statistics}]
DBCC SHOW_STATISTICS (
    'dbo.MRI_Scans',
    '_WA_Sys_00000004_4CA06362'
);
GO
\end{lstlisting}

The statistics metadata and density information, together with selected histogram entries, are summarized in Tables~\ref{tab:q2-statistics} and \ref{tab:q2_seriesdescription_histogram}, respectively.

\begin{table}[htbp]
\centering
\caption{Statistics Metadata and Density Information for \texttt{MRI\_Scans.SeriesDescription}}
\label{tab:q2-statistics}
\begin{tabular}{lr}
\toprule
Property & Value \\
\midrule
Rows & 48,345 \\
Rows sampled & 32,510 \\
Sampling percentage & 67.25\% \\
Steps & 199 \\
Density & 0.08955224 \\
All density & 0.003861004 \\
Average key length & 68 \\
String index & YES \\
Unfiltered rows & 48,345 \\
Persisted sample percent & 0 \\
\bottomrule
\end{tabular}
\end{table}

\begin{table}[htbp]
\centering
\small
\caption{Histogram Entries for \texttt{MRI\_Scans.SeriesDescription}}
\label{tab:q2_seriesdescription_histogram}
\begin{tabular}{|r|r|r|r|r|}
\hline
RANGE\_HI\_KEY & RANGE\_ROWS & EQ\_ROWS &
DISTINCT\_RANGE\_ROWS & AVG\_RANGE\_ROWS \\
\hline
T2\_TSE\_SAG\_384\_0002  & 0        & 7,256 & 0   & 1 \\
\hline
T2\_TSE\_TRA\_384\_0004  & 5,645.2701 & 7,145 & 1 & 5,645.2701 \\
\hline
T1\_TSE\_TRA\_0005       & 13,404.7807 & 7,049 & 2 & 3,829.615 \\
\hline
T1\_TSE\_SAG\_320\_0003  & 20,765.4104 & 6,878 & 3 & 2,115.4087 \\
\hline
\multicolumn{5}{|c|}{$\vdots$} \\
\hline
\multicolumn{2}{|r|}{last step} & 1 & 1 & 1 \\
\hline
\end{tabular}
\end{table}

The column holds 289 distinct values across 199 histogram steps, with a heavily skewed distribution: the first step alone (\texttt{T2\_TSE\_SAG\_384\_0002}) holds 7,256 rows, and the four most frequent series codes together cover 28,328 rows, approximately 58.6\% of the table. The tail of the histogram consists of many steps with \texttt{EQ\_ROWS = 1}, which the optimizer samples with reduced precision.

\subsection{Runtime and I/O Observations}

The \texttt{MRI\_Scans} scan performs the same 1,880 logical reads as Q1, because the scan visits the identical 1,876 leaf pages plus overhead. The difference is entirely in the result set: Q2 returns 15 rows instead of Q1's 65. The unsargable \texttt{LIKE} predicate therefore reduces the number of returned rows without reducing the I/O footprint, which confirms that any improvement must come from changing the access path (an index on the join column) rather than from the filter itself.

With a warm buffer pool, the captured run reports 0 physical reads for \texttt{MRI\_Scans} and 2 for \texttt{Patients}, with 15 ms of CPU time and 27 ms of elapsed time. The Q1 capture showed the same I/O footprint with an anomalous elapsed time of approximately 76.3 seconds dominated by a \texttt{SLEEP\_BPOOL\_STEAL} wait, while Q2 completes in under 30 ms. The I/O behaviour of the scan is thus stable across the two queries, and the observed elapsed time is strongly dependent on buffer pool state.

As a cross-reference, the missing-index hint for Q2 matches the covering index \texttt{IX\_MRIScans\_PatientID\_Covering} introduced by Role~2. With that index in place, the same query is served by an Index Seek that reads only the 65 rows belonging to patient \texttt{'0001'}, applying the \texttt{LIKE} predicate as a residual filter on the index rows instead of on the entire table.

% =============================================
% Q3: Nearest-Neighbor Feature Search
% =============================================

\section{Q3: Nearest-Neighbor Feature Search}

\subsection{Query Characteristics}

Q3 retrieves the feature vector for target scan \texttt{ScanID = 1}, calculates the Euclidean distance between that vector and the other feature vectors, joins the feature records with MRI scan metadata, and returns the three records with the smallest distances.

The distance expression uses four feature dimensions:

\[
d=\sqrt{
(f_1-T_1)^2+
(f_2-T_2)^2+
(f_3-T_3)^2+
(f_4-T_4)^2
},
\]

where $f_1,\ldots,f_4$ are the candidate feature values and $T_1,\ldots,T_4$ are the target feature values.

\subsection{Execution Plan and Join Algorithm}

Figure~\ref{fig:q3-plan} presents the baseline actual execution plan for the main nearest-neighbor query.

\begin{figure}[htbp]
\centering
\planfig{image/role_3/q3_execution_plan.png}{0.95\textwidth}{Q3 baseline parallel actual execution plan}
\caption{Baseline actual execution plan for Q3}
\label{fig:q3-plan}
\end{figure}

Unlike Q1 and Q2, the baseline plan uses a Hash Match (Inner Join) to combine \texttt{Image\_Features} and \texttt{MRI\_Scans}. Both tables are accessed through Clustered Index Scans. The results are repartitioned by \texttt{ScanID} before being processed by the parallel Hash Match operator.

The Hash Match operator estimates 43,475.4 output rows and produces 48,344 actual rows. Its Q-error is:

\[
Q\text{-error}=
\max\left(
\frac{48,344}{43,475.4},
\frac{43,475.4}{48,344}
\right)
\approx 1.11.
\]

The optimizer therefore underestimates the join output by approximately 10.1\%. Compared with Q1 and Q2, Q3 joins substantially larger inputs, making Hash Match an appropriate join strategy for this workload.

The relevant scan and join properties are summarized in Table~\ref{tab:q3-cardinality}.

\begin{table}[htbp]
\centering
\caption{Q3 scan and join cardinality}
\label{tab:q3-cardinality}
\begin{tabular}{lrr}
\toprule
Operator & Estimated rows & Actual rows \\
\midrule
\texttt{Image\_Features} scan & -- & 48,345 read; 48,344 returned \\
\texttt{MRI\_Scans} scan & -- & 48,345 read; 48,345 returned \\
Hash Match & 43,475.4 & 48,344 \\
\bottomrule
\end{tabular}
\end{table}

The \texttt{Image\_Features} scan applies \texttt{ScanID <> @TargetScanID} as a residual predicate, returning 48,344 rows. The \texttt{MRI\_Scans} scan reads all 48,345 rows.

\subsection{Parallel Execution}

The main Q3 query uses a degree of parallelism of eight. Figure~\ref{fig:q3-parallel} highlights the parallelism operators and data flow surrounding the Hash Match operator.

\newpage

\begin{figure}[htbp]
\centering
\planfig{image/role_3/q3_parallel_operators.png}{0.5\textwidth}{Q3 Repartition Streams, Hash Match, and Gather Streams}
\caption{Parallel data flow in the Q3 execution plan}
\label{fig:q3-parallel}
\end{figure}

The plan uses Repartition Streams to distribute rows by \texttt{ScanID} before the join and Gather Streams to collect the sorted results. The per-thread counters for the Hash Match operator sum to 48,344 rows across the eight active workers.

The Top-N Sort performs local sorting with \texttt{Top Rows = 3} on each worker. This produces up to 24 local candidates before the results are gathered. A final Top operator then returns the three globally closest records.

The separate statement that retrieves the target feature vector has a degree of parallelism of one. Therefore, the parallel execution described here applies to the main nearest-neighbor query rather than the preliminary lookup.

\subsection{Statistics Used for Cardinality Estimation}

The following query retrieves the statistics information for \texttt{Image\_Features.ScanID}, including its header, density vector, and histogram.

\begin{lstlisting}[caption={Retrieving Q3 statistics for Image\_Features.ScanID},
                   label={lst:q3-statistics}]
DBCC SHOW_STATISTICS (
    'dbo.Image_Features',
    'UQ__Image_Fe__63B32660DF05EFC2'
);
GO
\end{lstlisting}

The header of the statistics object \texttt{UQ\_\_Image\_Fe\_\_63B32660DF05EFC2}, associated with \texttt{Image\_Features.ScanID}, is summarized in Table~\ref{tab:q3-statistics}.

\begin{table}[htbp]
\centering
\caption{Statistics Metadata and Density Information for \texttt{Image\_Features.ScanID}}
\label{tab:q3-statistics}
\begin{tabular}{lr}
\toprule
Property & Value \\
\midrule
Rows & 48,345 \\
Rows sampled & 18,900 \\
Sampling percentage & 39.09\% \\
Steps & 175 \\
Density & 2.377895E-05 \\
Average key length & 8 bytes \\
String index & NO \\
Unfiltered rows & 48,345 \\
Persisted sample percent & 0 \\
\bottomrule
\end{tabular}
\end{table}

The histogram for \texttt{ScanID} contains 175 steps, with selected entries summarized in Table~\ref{tab:q3_scanid_histogram}. Each step describes the distribution of values up to its corresponding range boundary.

\begin{table}[htbp]
\centering
\small
\caption{Histogram Entries for \texttt{Image\_Features.ScanID}}
\label{tab:q3_scanid_histogram}
\begin{tabular}{|r|r|r|r|r|}
\hline
RANGE\_HI\_KEY & RANGE\_ROWS & EQ\_ROWS &
DISTINCT\_RANGE\_ROWS & AVG\_RANGE\_ROWS \\
\hline
40    & 0        & 1 & 0   & 1 \\
\hline
265   & 241.8147 & 1 & 224 & 1.07953 \\
\hline
655   & 264.9672 & 1 & 265 & 1.001525 \\
\hline
1070  & 162.0673 & 1 & 162 & 1 \\
\hline
1445  & 326.7071 & 1 & 326 & 1.002637 \\
\hline
\multicolumn{5}{|c|}{$\vdots$} \\
\hline
47322 & 326.7071 & 1 & 205 & 1.593693 \\
\hline
47580 & 326.7071 & 1 & 257 & 1.271234 \\
\hline
48020 & 326.7071 & 1 & 326 & 1.002637 \\
\hline
48296 & 272.6846 & 1 & 272 & 1.001691 \\
\hline
48297 & 0        & 1 & 0   & 1 \\
\hline
\end{tabular}
\end{table}

Although these statistics describe the indexed join column, the histogram alone does not explain the optimizer's complete cardinality estimate for the Q3 join. The execution plan also uses statistics from the other participating tables and columns. Therefore, the individual contribution of each statistics object to the estimated join cardinality cannot be determined from the available evidence.

\subsection{Compute Scalar and Vector Distance Calculation}

The Compute Scalar operator defines the expression \texttt{Expr1002}, which evaluates the four-dimensional Euclidean distance. The expression performs subtraction and squaring for each dimension, adds the four squared differences, and applies the square-root function. The resulting scalar value is used as the sort key for ranking candidate scans.

Figure~\ref{fig:q3-compute-scalar} presents the Compute Scalar operator and its expression properties.

\begin{figure}[htbp]
\centering
\planfig{image/role_3/q3_compute_scalar.png}{0.5\textwidth}{Q3 Compute Scalar expression properties}
\caption{Compute Scalar expression for the Euclidean distance}
\label{fig:q3-compute-scalar}
\end{figure}

The inspected plan exposes the defined expression and an estimated row count of 43,475.4. In the latest execution-plan capture, the Compute Scalar operator does not expose separate actual-row or execution counters. Consequently, no standalone actual-row count or runtime should be inferred from its properties.

The expression is evaluated as part of the query's processing pipeline before the Top-N Sort uses the calculated distances to select candidate records. This makes the scalar calculation an important part of Q3's workload, although its individual CPU cost cannot be isolated from the available operator properties.

\subsection{Top-N Sort and Result Selection}

The plan contains a Top-N Sort ordered by \texttt{Expr1002} in ascending order, with \texttt{Top Rows = 3}. Figure~\ref{fig:q3-sort} shows the sorting operator and its relevant properties.

\begin{figure}[htbp]
\centering
\planfig{image/role_3/q3_topn_sort.png}{0.5\textwidth}{Q3 Top-N Sort operator properties}
\caption{Top-N Sort used to rank candidate scans by Euclidean distance}
\label{fig:q3-sort}
\end{figure}

The parallel workers perform local Top-N sorting, retaining up to three candidates each. Gather Streams combines these local results, after which the final Top operator returns the three records with the smallest global distances.

This approach avoids sorting every candidate into a complete ordered result set. Instead, the execution plan retains a limited number of candidates from each worker before producing the final result.

\subsection{Reference-Vector Lookup}

Before executing the main nearest-neighbor query, SQL Server retrieves the four feature values for the target scan. This statement uses a Nested Loops plan consisting of an Index Seek on the unique nonclustered index associated with \texttt{Image\_Features.ScanID}, followed by a Clustered Index Seek on \texttt{PK\_Image\_Features} using the retrieved \texttt{FeatureID}.

Figure~\ref{fig:q3-lookup} shows the execution plan for this preliminary lookup.

\begin{figure}[htbp]
\centering
\planfig{image/role_3/q3_reference_lookup.png}{0.5\textwidth}{Q3 reference-vector lookup execution plan}
\caption{Execution plan for retrieving the target feature vector}
\label{fig:q3-lookup}
\end{figure}

Both seek operations return one row. The statement completes in approximately 3 milliseconds, with approximately 3 milliseconds of CPU time. Its execution is therefore negligible compared with the main nearest-neighbor query.

\subsection{Runtime, I/O, and Memory Observations}

The latest execution-plan capture reports an elapsed time of approximately 294 milliseconds and a CPU time of 751 milliseconds for the main query. CPU time exceeds elapsed time because the query executes in parallel, with CPU time accumulated across multiple worker threads.

The two Clustered Index Scans each read 48,345 rows. Their aggregate logical reads, calculated by summing the per-thread counters, are approximately 3,909 for \texttt{Image\_Features} and 1,972 for \texttt{MRI\_Scans}. The coordinator thread reports two physical reads for each table, while the worker-level physical reads are zero. The corresponding read-ahead reads are approximately 3,720 for \texttt{Image\_Features} and 1,874 for \texttt{MRI\_Scans}.

The main query receives a memory grant of 27,648 KB and reports a maximum memory usage of 8,128 KB. The plan also records a memory-grant wait of approximately 106 milliseconds. Thus, the granted memory exceeded the maximum amount used during this execution.

The wait statistics report \texttt{LATCH\_EX} as the largest wait, at approximately 1,517 milliseconds, followed by \texttt{CXSYNC\_PORT} at 327 milliseconds and \texttt{CXPACKET} at 168 milliseconds. These wait times are accumulated across execution threads and can therefore exceed the query's elapsed time. The execution plan does not establish the specific cause of the reported latch waits.

The latest capture also differs from the earlier observation, which reported approximately 78 milliseconds of elapsed time and 303 milliseconds of CPU time. Although both captures reused cached execution plans, their runtime measurements differ. The latest capture is used here as the current runtime observation; the difference indicates that execution time can vary between runs and should be evaluated using repeated measurements under consistent conditions.

\subsection{Hash Match versus Nested Loops for Q3}

The baseline Q3 execution plan uses a parallel Hash Match
(Inner Join) with a degree of parallelism of 8. To examine
the optimizer's join algorithm selection, an alternative
plan was generated using \texttt{OPTION (LOOP JOIN)}.
Table~\ref{tab:q3-join-comparison} compares their
estimated execution plan properties.

\begin{table}[htbp]
\centering
\caption{Q3: Hash Match versus Nested Loops}
\label{tab:q3-join-comparison}
\begin{tabular}{lrr}
\toprule
Metric & Hash Match & Nested Loops \\
\midrule
Estimated subtree cost & 11.7389 & 22.4434 \\
Join operator cost & 0.21353 (2\%) & 0.02569 (0\%) \\
Sort operator cost & -- & 6.8935 (31\%) \\
Memory grant (KB) & 27,648 & 1,024 \\
Degree of parallelism & 8 & 8 \\
Estimated rows & 43,475.4 & 43,475.4 \\
\bottomrule
\end{tabular}
\end{table}

The Hash Match plan has an estimated subtree cost of 11.7389, compared with 22.4434 for the forced Nested Loops plan. The Nested Loops alternative therefore has approximately 91.2\% higher estimated cost.

Although the Nested Loops operator has a lower estimated cost of 0.02569 than the Hash Match operator at 0.21353, its complete execution plan is more expensive which contains a Sort operator with an estimated cost of 6.8935, accounting for 31\% of the plan cost.

Both plans estimate 43,475.4 output rows. The baseline Hash Match plan produces 48,344 actual rows, resulting in a Q-error of approximately 1.112. The join algorithm change therefore does not alter the estimated cardinality.

Overall, the comparison supports the optimizer's choice of Hash Match because it has the lower estimated subtree cost. The Nested Loops plan's additional sorting requirement is an important contributor to its higher estimated cost.